Verified translation has two well-studied extremes: prove the translator
once (certified compilation), or validate each run of one translator
(translation validation). Both treat a single translation in isolation.
We study translations as a \emph{graph} --- many source languages,
several reasoning targets, multiple independently built routes --- where
the honest answer to ``is this translation correct?'' differs from edge
to edge. We present a calculus of \emph{fidelity-graded translations}:
pairs of languages close commuting squares that are checkable per
program and compose by pasting; declared fidelity grades compose by
weakest link, are \emph{re-established} per run by inline checking, and
are \emph{exceeded} by agreement between independently derived routes;
and an end-to-end theorem isolates a fundamental asymmetry ---
witness-carrying answers are self-certifying at the source, while
universal answers are where grades, branches, and certificates earn
their cost. The compositional core is mechanized in Lean~4.

The calculus is implemented in \hg, a platform of 13 languages and 13
pairs converging on two reasoning hubs, built as a \emph{two-directional
experiment in LLM-generated correctness}: independent LLM agents wrote
every pair, largely unsupervised, with the architecture's cross-checks
as the only semantic gate, and the platform's intended \emph{player} is
itself an LLM, handed deterministic, graded, checked moves. All code,
and most of this paper, is LLM-generated; the human contribution is the
architecture. We report per-construct coverage conjoined with fidelity,
branch agreement along dual RISC-V and AArch64 routes, source-level
witness replay, certified unreachability re-validated by a formally
verified checker, and the defects the architecture caught --- including
one in that checker's own adapter.
